\section{Confirmation and validation}
\label{sec:validation}

The empirical structure must not depend on one estimator, on the soft priors, or on the pooled cohort. Three
analyses address these in turn. Because the marginalized model and FIML share one objective
(Section~\ref{sec:estimation}), and a classical FIML on the full high-missingness backbone is intractable and
unreliable, the confirmation is performed \emph{in-engine}.

\subsection{Prior-free refit: is the structure earned, or manufactured by the priors?}
The continuous backbone (S2) is re-fit at full $N$ with \emph{flat} loading priors (identification constraints
only --- home cells stay sign-oriented, signed cells centred at $0$). A flat-prior MAP equals the MLE equals
FIML, so loadings and $\Phi$ that match the soft-prior fit show the structure is data-driven.

\begin{keyresult}[The map is not a Bayesian-prior artefact]
The prior-free full-$N$ refit reproduces the soft-prior fit \emph{exactly}: per-factor Tucker congruence
$\varphi(\text{soft},\text{flat})=1.00$ for every factor (overall, cognition, metabolic, inflammatory, sleep),
$\max|\Delta\Phi_{\text{off-diag}}|=0.000$, flat-fit $\widehat{R}=1.00$.
\end{keyresult}

\subsection{Absolute fit: posterior-predictive residual correlations}
Comparing each posterior draw's model-implied correlation to the observed pairwise-complete correlation gives a
Bayesian SRMR and a residual-correlation matrix --- absolute fit without the $\chi^2$ asymptotics that fail
under heavy missingness and non-normal indicators. The Bayesian SRMR is
$\mathbf{0.074}\ [0.073,0.075]$ (conventional good fit $<0.08$), and the residual misfit is confined to
\emph{repeated-measure clusters}: chol--LDL ($0.75$), supine/standing heart rate ($0.69$--$0.75$),
supine/standing blood pressure ($0.58$--$0.61$), WAIS subtests ($0.57$), AST--ALT ($0.55$). These are
method-factor artefacts (the same quantity measured twice), not structural misfit.

\subsection{Incremental fit: WAIC model comparison}
On a common $N=6{,}000$ random cohort-balanced sub-sample, the bifactor is compared with a unidimensional null
($G$ only) and a correlated-factors alternative (no general loadings on the specific-domain items).

\begin{table}[h]\centering\small\sffamily
\caption{WAIC comparison (lower is better). The bifactor structure is decisively preferred.}
\label{tab:waic}
\begin{tabular}{L{6.0cm} R{2.6cm} R{2.4cm} R{2.4cm}}
\toprule
\hdr{model} & \hdr{WAIC} & \hdr{$\Delta$WAIC} & \hdr{SE($\Delta$)}\\
\midrule
\rowcolor{facebg} bifactor ($G$ + specifics) & 702{,}633 & 0 & 0\\
correlated-factors (no $G$) & 705{,}352 & 2{,}720 & 78.9\\
\rowcolor{facebg} unidimensional ($G$ only) & 756{,}043 & 53{,}410 & 426.8\\
\bottomrule
\end{tabular}
\end{table}

Together: the continuous backbone is \textbf{estimator- and prior-robust} --- absolute fit is acceptable,
the structure reproduces under flat priors, and WAIC supports the bifactor over both simpler alternatives.
Classical CFI/TLI/RMSEA are not reported (convention, not new evidence, and statistically weak here); they are
available via \texttt{lavaan} on request.

\subsection{Measurement invariance across BP/SZ/DR}
Does each loading hold per cohort, or is the map partly cohort-driven? This is the central remaining validity
check. Classical multi-group SEM hits the same wall as FIML, so invariance is tested in-engine: the backbone is
fit \emph{per cohort} (z-scored within cohort) on a small cohort-balanced sub-sample over three seeds, and
metric invariance is read as Tucker congruence $\varphi$ of the primary loadings per factor per cohort-pair.
Simple structure (not bifactor) is used, as is conventional for multi-group invariance and because the
per-cohort bifactor $G$ is multimodal in SZ (no FAST anchor). All nine fits converged.

\begin{table}[h]\centering\small\sffamily
\caption{Metric invariance: Tucker congruence $\varphi$ (mean over seeds) per factor per cohort-pair
($\varphi\ge 0.95$ invariant; $\ge 0.85$ partial). $12/15$ comparisons fully invariant; $13/15$ at least
partial.}
\label{tab:invariance}
\begin{tabular}{L{3.6cm} R{1.7cm} R{1.7cm} R{1.7cm} L{3.6cm}}
\toprule
\hdr{factor} & \hdr{BP--SZ} & \hdr{BP--DR} & \hdr{SZ--DR} & \hdr{verdict}\\
\midrule
\rowcolor{facebg} overall severity & 0.92 & 0.99 & 0.97 & invariant (BP--SZ partial)\\
cognition & 0.99 & 0.96 & 0.97 & invariant\\
\rowcolor{facebg} metabolic & 0.99 & 0.97 & 0.99 & invariant\\
inflammatory & 0.99 & \textbf{0.71} & \textbf{0.75} & \textbf{non-invariant in DR}\\
\rowcolor{facebg} sleep & 1.00 & 0.99 & 0.99 & invariant\\
\bottomrule
\end{tabular}
\end{table}

\begin{caveatbox}[A documented partial-invariance caveat for DR inflammatory scores]
The dimensional structure is \emph{largely invariant} --- cognition, metabolic, and sleep recover congruently
in every cohort; $G$ is invariant except BP--SZ (partial, $\varphi=0.92$; few anchors, no FAST in SZ). The one
genuine exception is \emph{inflammatory in DR}: neutrophils load $\approx 0$ in DR ($0.07$ vs $0.88$ in BP/SZ)
while eosinophils load higher ($0.59$ vs $0.23$). DR's inflammatory axis is compositionally different, so DR
inflammatory scores carry a documented partial-invariance caveat rather than a hidden bias. Cohort-modular
dimensions (anhedonia; heterogeneously-measured suicidality) are declared modular, not claimed invariant.
\end{caveatbox}
